% ECONOMICS_PROBLEM: {"id": "R23-380", "title": "Mechanisms of neighborhood effects", "jel_code": "R23", "source_id": "OP-0466"}
% !TeX program = xelatex
\documentclass[11pt,a4paper]{article}
\usepackage[margin=23mm]{geometry}
\usepackage{fontspec}
\setmainfont{Latin Modern Roman}
\usepackage{amsmath,amssymb,mathtools}
\usepackage{xcolor,enumitem,booktabs,longtable,array,xurl,hyperref}
\definecolor{muted}{HTML}{53636C}
\hypersetup{colorlinks=true,urlcolor=blue,linkcolor=blue}
\setlength{\parindent}{0pt}
\setlength{\parskip}{5pt}
\newcommand{\R}{\mathbb R}
\newcommand{\E}{\mathbb E}
\newcommand{\N}{\mathbb N}
\newcommand{\ind}{\mathbf 1}
\newcommand{\Var}{\operatorname{Var}}
\newcommand{\Cov}{\operatorname{Cov}}
\newcommand{\supp}{\operatorname{supp}}
\DeclareMathOperator*{\argmin}{arg\,min}
\DeclareMathOperator*{\argmax}{arg\,max}
\newcommand{\entrymeta}[3]{{\sffamily\small\color{muted}\textbf{\texttt{#1}}\quad|\quad #2\par #3\par}}
\newcommand{\Problemref}[1]{\texttt{#1}}
\newcommand{\JELref}[2]{\texttt{#2} (JEL \texttt{#1})}
\begin{document}
\section*{Mechanisms of neighborhood effects}
\textbf{Problem ID:} \texttt{R23-380}\quad \textbf{JEL:} \texttt{R23}\par
% BEGIN ECONOMICS STATEMENT
\label{problem:OP-0466}
{\sffamily\small\color{muted}Original subject group: Urban, housing and spatial economics\par}
\label{definitions:problem:30007748}
\textbf{Audit classification:} Published research agenda. Added scale-correct mechanism derivatives, feasible component interventions and path-dependent finite-change attribution.
\subsection*{Definitions and precise research target}
\textit{Editorial formalization.} Consider children aged six through twelve whose families receive randomized offers to move within one metropolitan area. Let neighborhood exposure be a vector \(x=(s,p,v,e)\), where \(s\) measures school quality, \(p\) peer composition, \(v\) exposure to violence, and \(e\) environmental quality using independently defined scales. Let \(Y_i(x)\) be adult earnings under a sustained six-year exposure profile. For baseline vector \(x_0\) and a feasible intervention changing only component \(j\), write \(x_0+\delta e_j\), where \(e_j\) is the corresponding unit vector. Define
\[\theta_j=\left.\frac{\partial E[Y_i(x_0+\delta e_j)]}{\partial\delta}\right|_{\delta=0}.\]
Require a feasible neighborhood of the component intervention and existence of the mean derivative; use one-sided contrasts at boundaries. Identify the relative contributions of these components and their interactions to the total move effect, using independent interventions or explicit valid instruments. A voucher lottery shifts several components together and ordinarily identifies a bundled effect rather than each \(\theta_j\). Specify exposure timing, household responses, and exclusion restrictions, and report bounds if component variation is insufficient. Evaluate whether results transport beyond the stated age range and metropolitan setting.

Each exposure coordinate has a fixed measurement scale and an implementable intervention holding other coordinates fixed. Let \(m(x)=E[Y_i(x)]\) for a sustained six-year exposure profile with fixed baseline child weights; require differentiability on a feasible open neighborhood of \(x_0\). Then \(\theta_j=\partial_jm(x_0)\) is expressed in earnings units per unit of coordinate \(j\); raw derivatives in different units are not directly comparable. For a finite bundled change \(x_1-x_0\), define an attribution path \(x(u)\) and contributions \(c_j=\int_0^1\partial_jm(x(u))x_j'(u)du\), so \(\sum_jc_j=m(x_1)-m(x_0)\) when the chain rule applies. Path choice matters with interactions. A single move lottery does not separately identify four derivatives without rank, exclusion and support restrictions. All expectations use a stated baseline population and probability law. Identification means every admissible data-generating model with the same observed-data law yields the same displayed target; otherwise report the identified set. Exact equations, horizons and assumptions are editorial, rather than source-given canonical conjectures.
\subsection*{Variants and assumption changes}
\begin{enumerate}
\item \textbf{Editorial finite-difference variant.} Replace infeasible coordinate derivatives by implementable bundles \(x_0,x_1\) and estimate \(m(x_1)-m(x_0)\); this scalar contrast does not separate mechanisms.
\item \textbf{Editorial timing variant.} Compare equal cumulative exposure at ages 6–12 versus 12–18, defining complete age profiles. Duration alone no longer summarizes treatment.
\end{enumerate}
\subsection*{References and source audit}
\textbf{1.} Explicit agenda on mechanism weights, scaling mobility programs, and incumbent residents of revitalized areas. Locator: Journal of Economic Perspectives 35(4), 197–222; Discussion and Conclusion, pp.216–217. Full text or primary publication page inspected. URL: \url{https://lkatz.scholars.harvard.edu/sites/g/files/omnuum5961/files/lkatz/files/chyn_katz_jep_2021_all.pdf}.

\textbf{2.} Opportunity-market review exists; older review supplies explicit agenda, not this abstract. Locator: NBER Working Paper 35419, July 2026; abstract. Primary indexed metadata/abstract inspected; direct record failed. URL: \url{https://www.nber.org/papers/w35419}.
\begin{enumerate}\raggedright
\item\hypertarget{definitions:source:30007748:1}{} Eric Chyn; Lawrence F. Katz. 2021. \emph{Neighborhoods Matter: Assessing the Evidence for Place Effects}. Journal of Economic Perspectives 35(4), 197–222; Discussion and Conclusion, pp.216–217.
\url{https://lkatz.scholars.harvard.edu/sites/g/files/omnuum5961/files/lkatz/files/chyn_katz_jep_2021_all.pdf}.
DOI: \url{https://doi.org/10.1257/jep.35.4.197}.
\item\hypertarget{definitions:source:30007748:2}{} Lawrence F. Katz. 2026. \emph{Neighborhood Effects and Missing Markets for Opportunity}. NBER Working Paper 35419, July 2026; abstract.
\url{https://www.nber.org/papers/w35419}.
DOI: \url{https://doi.org/10.3386/w35419}.
\end{enumerate}

\subsection*{Common conventions: Source mathematical conventions}

These conventions apply to each entry unless explicitly replaced.

\textbf{Models and identification.} A primitive model \(m\) specifies all probability laws, feasible choices, information, equilibrium and measurement rules used by its target. The admissible class \(\mathcal M\) is fixed independently of outcomes. If \(P_O(m)\) is its observable probability law and \(\tau(m)\) the target, its identified set at observable law \(P\) is
\[
 \mathcal I_\tau(P)=\{\tau(m):m\in\mathcal M,\ P_O(m)=P\}.
\]
Point identification means that this set is a singleton. An empty set signals incompatibility of data and assumptions. Coordinate bounds need not describe the joint set. Bounds are sharp only if every retained target is attainable or arbitrarily approachable by compatible models. Finite bounds require explicit restrictions on unobserved quantities; unrestricted confounding, hidden wealth, extrapolation or arbitrary equilibrium selection can yield uninformative sets.

\textbf{Causal interventions.} A potential outcome \(Y_i(p)\) is the outcome under a complete policy regime \(p\), including timing, financing, information and market spillovers. The target population and units are fixed before treatment. With interference, use the complete assignment vector or a justified exposure mapping; individual no-interference notation alone cannot identify an economy-wide intervention. A difference of mean potential outcomes does not require the cross-world joint distribution. Conditioning on post-treatment migration, survival, enrollment or employment changes the estimand and needs separate justification.

\textbf{Statistical statements.} An estimator, confidence set, forecast or test specifies a sampling law, model class and loss/coverage criterion. Uniform \((1-\alpha)\)-coverage means \(\inf_{m\in\mathcal M}\Pr_m\{\tau(m)\in C_n\}\geq1-\alpha\), or an explicitly asymptotic version. The whole parameter space trivially has coverage, so informativeness requires a stated width, risk or power objective. A forecasting improvement is relative to a named benchmark, information set and evaluation population.

\textbf{Equilibrium and optimization.} An equilibrium comprises feasible plans, optimality, beliefs where needed, budget constraints and all market/resource clearing. The primitives or a stated benchmark define each correspondence \(\mathcal E(p;m)\). Nonemptiness is assumed or proved, never inferred just from compact policy sets. A selected-equilibrium target uses a declared selection rule; a robust target quantifies over all permitted equilibria. Use suprema or infima unless attainment is established. An \(\varepsilon\)-optimal policy is feasible and within \(\varepsilon\) of the finite optimum value. Report an empty feasible set separately.

\textbf{Welfare and accounting.} Normative utility, interpersonal normalization, social weights, numeraire, discounting and the use of public receipts are primitives. Behavioral utility need not coincide with the welfare criterion. Household welfare includes ownership income and rents once; adding profits again to compensating variation generally double-counts them. A monetary equivalent is defined by an explicit transfer and price convention, with monotonicity and range conditions. Average effects, mechanism contributions, transfers and welfare losses are different targets. Infinite sums require convergence and dynamic feasible plans require terminal or no-Ponzi conditions.

\textbf{Source versions.} A working-paper date, revision date and journal publication year are distinguished. A DOI for a journal version is not automatically the DOI of an earlier manuscript. Locators refer to the stated version and printed pages where specified. A metadata-only check does not verify a claim about a full-text conclusion. Every entry's source subsection narrows the provenance claim accordingly.
% END ECONOMICS STATEMENT
\end{document}
